\errorcontextlines=99999
\documentclass[a4paper,10pt,parskip=half,numbers=noenddot]{scrartcl}

\usepackage{hyperref}
\usepackage{xspace}
\usepackage{xcolor}
\usepackage{listings}
\usepackage{needspace}
\usepackage[most]{tcolorbox}
\usepackage{tikzquack}

\newcommand\pkg[1]{\textsf{#1}}
\newcommand\cs[1]{\texttt{\textbackslash#1}}

\newcommand{\pkgname}{\pkg{tikzquack}\xspace}


\lstdefinestyle{myexample}{
	language=[LaTeX]TeX,
	basicstyle=\ttfamily\small,
	keywordstyle=\color{blue!60!black},
	commentstyle=\color{green!60!black},
	tabsize=4
}

\newtcblisting{example}{
	listing and text,
	listing options={
		style=myexample
	},
	colback=white,
	colframe=black!75,
	colbacklower=white
}

\title{Il pacchetto \pkgname}
\author{Luca Robbiano\\\url{https://github.com/lr94/tikzquack}}
\date{v1.0, 2026/11/07}

\begin{document}
\maketitle

Un piccolo e semplice pacchetto \LaTeX\xspace per disegnare la famiglia dei paperi con TikZ :)

\section{Utilizzo}
	\Needspace{6\baselineskip}
	Il pacchetto \pkgname fornisce un comando \cs{papero}. Di default viene disegnato \mbox{Peppone}:
	\begin{example}
\begin{tikzpicture}
	\papero
\end{tikzpicture}
	\end{example}
	
	\Needspace{6\baselineskip}
	È possibile specificare quale papero disegnare e applicare le classiche opzioni TikZ per dimensione, posizione o altre trasformazioni:
	\begin{example}
\begin{tikzpicture}
	\papero[peppone]
	\papero[piumotta,scale=0.5,shift={(6,0)}]
	\papero[peppe,scale=0.8,shift={(7,0)}]
	\papero[otto,scale=0.5,shift={(17,0)}]
\end{tikzpicture}
	\end{example}
	
	\Needspace{6\baselineskip}
	Solo alcuni dei paperi sono definiti (quelli dell'esempio precedente). Però è possibile definire un nuovo papero col comando \cs{newpapero} specificando in HTML i suoi due colori:
	\begin{example}
\newpapero{piumina}{ffc0cb}{8aff88}
\begin{tikzpicture}
\papero[piumina]
\end{tikzpicture}
	\end{example}
	
	\Needspace{6\baselineskip}
	In alternativa è possibile specificare i due colori (non necessariamente in HTML, qui) e anche quelli di occhi e nastrino (che è anche possibile nascondere):
	\begin{example}
\begin{tikzpicture}
\papero[% senza indentazione, vedi "problemi noti"
primary=pink,%
secondary=green,%
ribboncolor=cyan,%
eyecolor=cyan%
]

% Piumotta con nastrino rosa
\papero[piumotta,ribboncolor=pink,shift={(4,0)}]

% Peppone senza nastrino
\papero[peppone,no ribbon,shift={(8,0)}]
\end{tikzpicture}
	\end{example}
	
	\Needspace{16\baselineskip}
	Oltre alla visuale principale (dall'alto) ne sono disponibili altre:
	\begin{example}
\begin{tikzpicture}
	\papero[peppone,top]
	\papero[peppone,front,shift={(3.2,0)}]
	\papero[piumotta,tophead,shift={(6,0)}]
	\papero[piumotta,lowerhead,shift={(9,0)}]
	% posa intenzionalmente meno "simmetrica"
	\papero[peppone,natural,shift={(4.7,-4)}]
\end{tikzpicture}
	\end{example}
	
\section{Problemi noti}
	Potrebbero esserci ancora alcuni problemi non del tutto risolti quando viene usato il comando \cs{newpapero} in un ambiente che crea un gruppo.\\
	Non è possibile spezzare le opzioni su più righe e indentarle. Se si spezzano le opzioni su più righe è necessario mettere un \% alla fine di ogni riga e fare iniziare la seguente all'inizio della riga stessa, senza spazi.
\end{document}